% Part: incompleteness
% Chapter: representability-in-q
% Section: sigma1-completeness

\documentclass[../../../include/open-logic-section]{subfiles}

\begin{document}

\olfileid{inc}{inp}{s1c}
\olsection{تمام \texorpdfstring{$\Sigma_1$}{Sigma-1}}

على الرغم من عدم تمام $\Th{Q}$ وامتداداتها المتسقة القابلة للمحورة،
رأينا أن $\Th{Q}$ تثبت حقائق أساسية كثيرة عن الأعداد الاسمية. ويمكن
في الواقع توسيع ذلك إلى حد بعيد. ولكي نفهم نطاق ما يمكن إثباته في
$\Th{Q}$، ندخل مفاهيم !!{formula}s من الأصناف $\Delta_0$ و
$\Sigma_1$ و$\Pi_1$. وعلى وجه التقريب، تكون !!{formula} من الصنف
$\Sigma_1$ إذا كانت من الصورة $\lexists[x][!B(x)]$، حيث تُبنى $!B$
باستعمال الروابط القضوية والمكممات المحدودة وحدها. وسنبيّن أنه إذا
كانت $!A$ !!{sentence} من الصنف $\Sigma_1$ وصادقة في $\Struct{N}$،
فإن $\Th{Q}\Proves !A$
(\olref{thm:sigma1-completeness}).

\begin{defn}
\ollabel{defn:bd-quant}
تكون \emph{!!{formula} وجودية محدودة} من الصورة
$\lexists[x][(x<t\land !A(x))]$، حيث $t$ حد كيفما كان، ونكتبها
اصطلاحًا $\bexists{x<t}{!A(x)}$.
%
وتكون \emph{!!{formula} كلية محدودة} من الصورة
$\lforall[x][(x<t\lif !A(x))]$، حيث $t$ حد كيفما كان، ونكتبها
اصطلاحًا $\bforall{x<t}{!A(x)}$.
\end{defn}

\begin{defn}
\ollabel{defn:delta0-sigma1-pi1-frm}
تكون !!{formula}~$!B$ من الصنف $\Delta_0$ إذا بُنيت من
!!{formula}s ذرية باستعمال الروابط القضوية والتكميم المحدود وحدهما.
%
وتكون !!{formula}~$!A$ من الصنف $\Sigma_1$ إذا
$!A\ident\lexists[x][!B(x)]$، حيث $!B$ من الصنف $\Delta_0$.
%
وتكون !!{formula}~$!A$ من الصنف $\Pi_1$ إذا
$!A\ident\lforall[x][!B(x)]$، حيث $!B$ من الصنف $\Delta_0$.
\end{defn}

\begin{lem}
\ollabel{lem:q-proves-clterm-id} افترض أن $t$ حد مغلق بحيث
$\Value{t}{N}=n$. عندئذ $\Th{Q}\Proves\eq[t][\num n]$.
\end{lem}

\begin{proof}
نثبت ذلك بالاستقراء على تعقيد~$t$. في الحالة الأساسية
$\Value{\Obj 0}{N}=0$، و$\Th{Q}\Proves\eq[\Obj 0][\num 0]$ لأن
$\num 0\ident\Obj 0$.
%
وفي حالة الاستقراء، لتكن $t_1$ و$t_2$ حدين بحيث
$\Value{t_1}{N}=n_1$ و$\Value{t_2}{N}=n_2$، و
$\Th{Q}\Proves\eq[t_1][\num{n_1}]$، و
$\Th{Q}\Proves\eq[t_2][\num{n_2}]$.

عندئذ $\Value{(t_1')}{N}=n_1+1$، ولدينا أن
$\Th{Q}\Proves\eq[t_1'][{\num n_1}']$ بقواعد الرتبة الأولى
للهوية مطبقةً على فرض الاستقراء وعلى !!{formula}
$\eq[\num{n_1}'][\num{n_1}']$؛ ومن ثم
$\Th{Q}\Proves\eq[t_1'][\num{n_1+1}]$ بحسب تعريف الأعداد الاسمية.

وفي حالة المجاميع لدينا
\[
      \Value{(t_1 + t_2)}{N}
    = \Value{t_1}{N} + \Value{t_2}{N}
    = n_1 + n_2.
\]
وبفرض الاستقراء وقواعد الهوية،
$\Th{Q}\Proves\eq[t_1+t_2][\num{n_1}+t_2]$، ثم
$\Th{Q}\Proves\eq[t_1+t_2][\num{n_1}+\num{n_2}]$
بتطبيق ثان لقواعد الهوية. وبحسب
\olref[inc][req][bre]{lem:q-proves-add}،
$\Th{Q}\Proves
\eq[\num{n_1}+\num{n_2}][\num{n_1+n_2}]$، ومن ثم
$\Th{Q}\Proves\eq[t_1+t_2][\num{n_1+n_2}]$.

وتصلح حجة مماثلة أيضًا لـ$\times$، باستعمال
\olref[inc][req][bre]{lem:q-proves-mult}.
%
ولما كان هذا يستنفد حدود الحساب المغلقة، نحصل على
$\Th{Q}\Proves\eq[t][\num n]$ لكل حد مغلق~$t$ بحيث
$\Value{t}{N}=n$.
\end{proof}

\begin{prob}
أثبت بالتفصيل الجزء المتعلق بـ$\times$ من
\olref{lem:q-proves-clterm-id}.
\end{prob}

\begin{lem}
\ollabel{lem:atomic-completeness}
افترض أن $t_1$ و$t_2$ حدان مغلقان. عندئذ
\begin{enumerate}
\item إذا كان $\Value{t_1}{N}=\Value{t_2}{N}$،
    فإن $\Th{Q}\Proves\eq[t_1][t_2]$.
\item إذا كان $\Value{t_1}{N}\neq\Value{t_2}{N}$،
    فإن $\Th{Q}\Proves\eq/[t_1][t_2]$.
\item إذا كان $\Value{t_1}{N}<\Value{t_2}{N}$،
    فإن $\Th{Q}\Proves t_1<t_2$.
\item إذا كان $\Value{t_2}{N}\leq\Value{t_1}{N}$،
    فإن $\Th{Q}\Proves\lnot(t_1<t_2)$.
\end{enumerate}
\end{lem}

\begin{proof}
إذا أُعطينا الحدين $t_1$ و$t_2$، ثبّتنا
$n=\Value{t_1}{N}$ و$m=\Value{t_2}{N}$.

افترض أن $!A\ident t_1=t_2$. بحسب
\olref{lem:q-proves-clterm-id}،
$\Th{Q}\Proves\eq[t_1][\num n]$ و
$\Th{Q}\Proves\eq[t_2][\num m]$. فإذا كان $n=m$، فإن
$\Th{Q}\Proves\eq[\num n][\num m]$، ومن ثم
$\Th{Q}\Proves\eq[t_1][t_2]$ بتعدي الهوية. وإذا كان $n\neq m$،
فإن $\Th{Q}\Proves\eq/[\num n][\num m]$، وبتعدي الهوية مرة أخرى
$\Th{Q}\Proves\eq/[t_1][t_2]$.

والآن لتكن $!A\ident t_1<t_2$. نعتمد في الحالتين على
البديهية~$!Q_8$، التي تقرر
$x<y\liff\lexists[z][\eq[z'+x][y]]$ لكل $x,y$.

افترض $\Sat{N}{t_1<t_2}$. عندئذ يوجد $k\in\Nat$ بحيث
$n+k+1=m$. وبحسب \olref{lem:q-proves-clterm-id}،
$\Th{Q}\Proves\eq[t_1][\num n]$ و
$\Th{Q}\Proves\eq[t_2][\num m]$، وبحسب الجزء الأول من هذه اللمة،
$\Th{Q}\Proves\eq[{\num k}'+\num n][\num m]$. ويترتب من تعدي
الهوية أن
$\Th{Q}\Proves\eq[{\num k}'+t_1][t_2]$، ومن ثم
$\Th{Q}\Proves\lexists[z][\eq[z'+t_1][t_2]]$. وبالاتجاه من اليمين
إلى اليسار من~$!Q_8$، $\Th{Q}\Proves t_1<t_2$.

وافترض بدلًا من ذلك $\Sat/{N}{t_1<t_2}$؛ أي $m\leq n$.
%
نعمل في~$\Th{Q}$ ونفترض $t_1<t_2$. بالاتجاه من اليسار إلى اليمين
من~$!Q_8$، يوجد~$z$ بحيث $\eq[z'+t_1][t_2]$. ولما كانت
$\Th{Q}\Proves\eq[t_1][\num n]$ و
$\Th{Q}\Proves\eq[t_2][\num m]$، نحصل على
$\eq[z'+\num n][\num m]$.
%
وباستقراء خارجي على~$m$ يستعمل~$!Q_5$،
$\eq[z'+\num{n-m}][\Obj 0]$. فإذا كان $m=n$، كانت
$\eq[z'][\Obj 0]$، وهو تناقض مع~$!Q_2$. وإذا كان $m<n$، كانت
$\eq[(z'+\num{n-m-1})'][\Obj 0]$ بحسب~$!Q_5$ مرة أخرى، وهو تناقض
مع~$!Q_2$. ومن ثم $\Th{Q}\Proves\lnot(t_1<t_2)$.
\end{proof}

\begin{lem}
\ollabel{lem:bounded-quant-equiv}
افترض أن $!A$ !!a{formula}، وأن $t$ حد مغلق، وأن
$k=\Value{t}{N}$. عندئذ
\begin{enumerate}
\item $\Th{Q}\Proves\bforall{x<t}{!A(x)}$ إذا وفقط إذا
    $\Th{Q}\Proves !A(\num 0)\land\dots\land !A(\num{k-1})$.
\item $\Th{Q}\Proves\bexists{x<t}{!A(x)}$ إذا وفقط إذا
    $\Th{Q}\Proves !A(\num 0)\lor\dots\lor !A(\num{k-1})$.
\end{enumerate}
\end{lem}

\begin{proof}
    نثبت حالة المكمم الكلي المحدود. فإذا كانت $\Value{t}{N}=0$،
    كان الطرف الأيسر من التكافؤ قابلًا للإثبات في~$\Th{Q}$، لأنه
    لا يوجد $x<\num 0$ بحسب
    \olref[inc][req][min]{lem:less-zero}. وبالمثل، يمكننا أن نأخذ
    الاقتران الخالي على أنه $\ltrue$ فحسب، وهي قابلة للإثبات أيضًا
    في~$\Th{Q}$.
    %
    ولذلك نفترض أن $\Value{t}{N}=k+1$ لعدد طبيعي ما~$k$. وبحسب
    \olref{lem:q-proves-clterm-id} يجوز أن نفترض أننا نعمل مع
    !!a{formula} من الصورة $\bforall{x<\num{k+1}}{!A(x)}$.
    
    افترض أن
    $\Th{Q}\Proves\bforall{x<\num{k+1}}{!A(x)}$، ولتكن
    $n\leq k$. ولما كانت
    $\Th{Q}\Proves\num n<\num{k+1}$ بحسب
    \olref{lem:atomic-completeness}، ينتج بالمنطق أن
    $\Th{Q}\Proves !A(\num n)$. وبتطبيق هذه الحقيقة $k+1$ مرة،
    مرة لكل $n\leq k$، نحصل على
    $\Th{Q}\Proves !A(\num 0)\land\dots\land !A(\num k)$ كما
    هو مطلوب.
    
    وفي الاتجاه الآخر، افترض أن
    $\Th{Q}\Proves !A(\num 0)\land\dots\land !A(\num k)$.
    وبالعمل في $\Th{Q}$، افترض $x<\num{k+1}$. بحسب
    \olref[inc][req][min]{lem:less-nsucc} لدينا
    $x=\num 0\lor\dots\lor x=\num k$، فينتج بالمنطق~$!A(x)$،
    ومن ثم الدعوى الكلية
    $\bforall{x<\num{k+1}}{!A(x)}$.
    
    وبرهان التكافؤ لـ!!{formula}s المكممة وجوديًا تكميمًا محدودًا
    مماثل.
\end{proof}

\begin{prob}
قدّم برهانًا مفصلًا للحالة الوجودية في
\olref{lem:bounded-quant-equiv}.
\end{prob}

\begin{lem}
\ollabel{lem:delta0-completeness}
إذا كانت $!A$ !!{sentence} من الصنف $\Delta_0$ وصادقة في
$\Struct{N}$، فإن $\Th{Q}\Proves !A$.
\end{lem}

\begin{proof}
نثبت ذلك بالاستقراء على تعقيد !!{formula}.
%
الحالة الأساسية معطاة بـ\olref{lem:atomic-completeness}، ولذلك ننتقل
إلى خطوة الاستقراء. وللتبسيط نقسم حالة النفي إلى حالات فرعية بحسب
بنية !!{formula} التي يطبق عليها النفي.

\begin{enumerate}
\item افترض أن $(!A\land !B)$ صادقة في $\Struct{N}$؛ ومن ثم تكون
$!A$ و$!B$ صادقتين في~$\Struct{N}$. وبفرض الاستقراء،
$\Th{Q}\Proves !A$ و$\Th{Q}\Proves !B$، ولذلك
$\Th{Q}\Proves(!A\land !B)$ بالمنطق.
%
\item افترض أن $\lnot(!A\land !B)$ صادقة في $\Struct{N}$؛ ومن ثم
تكون إما $\lnot !A$ وإما $\lnot !B$ صادقة في~$\Struct{N}$. ومن دون
إخلال بالعموم، افترض الأولى. بفرض الاستقراء
$\Th{Q}\Proves\lnot !A$، ومن ثم
$\Th{Q}\Proves\lnot(!A\land !B)$ بالمنطق.
%
\item افترض أن $(!A\lor !B)$ صادقة في $\Struct{N}$؛ فتكون إما
$!A$ صادقة في $\Struct{N}$ وإما $!B$ صادقة في $\Struct{N}$. ومن دون إخلال
بالعموم، افترض الأولى. بفرض الاستقراء $\Th{Q}\Proves !A$، ومن ثم
$\Th{Q}\Proves(!A\lor !B)$ بالمنطق.
%
\item افترض أن $\lnot(!A\lor !B)$ صادقة في $\Struct{N}$؛ فتكون
$\lnot !A$ و$\lnot !B$ صادقتين في $\Struct{N}$. عندئذ
$\Th{Q}\Proves\lnot !A$ و$\Th{Q}\Proves\lnot !B$ بفرض
الاستقراء. وبالتالي $\Th{Q}\Proves\lnot(!A\lor !B)$ بالمنطق.
%
\item افترض أن $\bforall{x<t}{!A(x)}$ صادقة في~$\Struct{N}$، حيث
$t$ حد مغلق و$k=\Value{t}{N}$. بفرض الاستقراء والمنطق، إذا كانت
$!A(\num n)$ صادقة في~$\Struct{N}$ لكل
$n<\Value{t}{N}$، فإن
$\Th{Q}\Proves !A(\num 0)\land\dots\land !A(\num{k-1})$.
وبحسب \olref{lem:bounded-quant-equiv} ينتج
$\Th{Q}\Proves\bforall{x<t}{!A(x)}$.
%
\item حالة المكمم الوجودي المحدود، حيث لدينا !!a{sentence} من الصورة
$\bexists{x<t}{!A(x)}$، مماثلة لحالة المكمم الكلي المحدود.
%
\item افترض أن $\lnot\bforall{x<t}{!A(x)}$ صادقة في
$\Struct{N}$، حيث $t$ حد مغلق. تكافئ هذه !!{sentence} الـ!!{sentence} الصيغة
$\bexists{x<t}{\lnot !A(x)}$، والتكافؤ قابل للاشتقاق في~$\Th{Q}$؛
فيجوز أن نطبق حجة المكممات الوجودية المحدودة.
%
\item وبالمثل، افترض أن $\lnot\bexists{x<t}!A(x)$ صادقة في
$\Struct{N}$، حيث $t$ حد مغلق. تكافئ هذه !!{sentence} في
$\Th{Q}$ الصيغة $\bforall{x<t}{\lnot !A(x)}$، فيجوز أن نطبق حجة
المكممات الكلية المحدودة.
%
\item أخيرًا، افترض أن $\lnot !A$ صادقة في $\Struct{N}$. لا يبقى
إلا حالتان: أن تكون $!A$ ذرية، وأن يكون
$\lnot !A\ident\lnot\lnot !B$ لبعض !!{sentence}~$!B$ من الصنف
$\Delta_0$. فإذا كانت $!A$ ذرية، فإن
$\Th{Q}\Proves\lnot !A$ بحسب \olref{lem:atomic-completeness}. وإذا
كان $\lnot !A\ident\lnot\lnot !B$، فهي تكافئ~$!B$ تكافؤًا قابلًا
للإثبات بالمنطق في~$\Th{Q}$، وتكون $!B$ صادقة في~$\Struct{N}$ لأن
$\lnot !A$ صادقة في~$\Struct{N}$. ومن ثم يعطينا فرض الاستقراء
$\Th{Q}\Proves\lnot !A$.
\end{enumerate}
\end{proof}

\begin{prob}
قدّم برهانًا مفصلًا للحالة الوجودية في
\olref{lem:delta0-completeness}.
\end{prob}

\begin{thm}
\ollabel{thm:sigma1-completeness}
إذا كانت $!A$ !!{sentence} من الصنف $\Sigma_1$ وصادقة في
$\Struct{N}$، فإن $\Th{Q}\Proves !A$.
\end{thm}

\begin{proof}
إذا كانت $\lexists{x}!A(x)$ !!{sentence} من الصنف $\Sigma_1$
وصادقة في~$\Struct{N}$، وُجد عدد طبيعي~$n$ وإسناد متغيرات~$s$ بحيث
$s(x)=n$ و$\Sat{N}{!A(x)}[s]$. وبحسب الحقائق القياسية عن علاقة
التحقيق ينتج $\Sat{N}{!A(\num n)}$. لكن $!A(\num n)$
!!{formula} من الصنف $\Delta_0$، ولذلك، بحسب
\olref{lem:delta0-completeness}، لدينا
$\Th{Q}\Proves !A(\num n)$، ومن ثم، بالمنطق، لدينا أيضًا
$\Th{Q}\Proves\lexists[x][!A(x)]$.
\end{proof}

\end{document}
